% =============================================================================
% preamble.tex — charte graphique du rapport (inspirée du Datachallenge_rapport)
% Moteur : lualatex (fontspec). Inclus via pandoc --include-in-header.
% =============================================================================

% ---- Couleurs de la charte ---------------------------------------------------
\usepackage{xcolor}
\definecolor{accent}{RGB}{38,63,140}     % indigo/bleu — titre + sections
\definecolor{accent2}{RGB}{95,80,160}    % violet — labels d'encadré
\definecolor{boxbg}{RGB}{238,238,248}    % lavande clair — fond d'encadré
\definecolor{boxframe}{RGB}{198,198,224}
\definecolor{shadecolor}{RGB}{245,246,248} % fond des blocs de code
\definecolor{rulecol}{RGB}{198,198,224}

% ---- Polices (lualatex + fontspec) ------------------------------------------
\usepackage{fontspec}   % défaut = Latin Modern (comme la charte de référence)
\setmonofont{DejaVu Sans Mono}[Scale=0.82]  % couvre → ─ │ ✓ des logs/code
\usepackage{microtype}

% ---- Symboles Unicode (emoji/flèches) hors verbatim -------------------------
\usepackage{pifont}
\usepackage{newunicodechar}
\newunicodechar{✅}{{\color{green!55!black}\ding{51}}}
\newunicodechar{❌}{{\color{red!72!black}\ding{55}}}
\newunicodechar{🛑}{{\color{red!72!black}\ding{54}}}
\newunicodechar{✓}{{\color{green!55!black}\ding{51}}}
\newunicodechar{✗}{{\color{red!72!black}\ding{55}}}
\newunicodechar{→}{\ensuremath{\rightarrow}}
\newunicodechar{←}{\ensuremath{\leftarrow}}
\newunicodechar{↑}{\ensuremath{\uparrow}}
\newunicodechar{🎯}{{\color{accent}\ding{110}}}
\newunicodechar{⚠}{{\color{orange!80!black}\ding{224}}}
\newunicodechar{≈}{\ensuremath{\approx}}
\newunicodechar{≥}{\ensuremath{\geq}}
\newunicodechar{≤}{\ensuremath{\leq}}
\newunicodechar{×}{\ensuremath{\times}}
\newunicodechar{↔}{\ensuremath{\leftrightarrow}}
\newunicodechar{«}{\guillemotleft}
\newunicodechar{»}{\guillemotright}

% ---- Titres de section : bleus, sans numéro LaTeX (numéros manuels conservés)
\usepackage{titlesec}
\setcounter{secnumdepth}{-1}   % on garde les numéros MANUELS dans le texte
\titleformat{\section}{\normalfont\Large\bfseries\color{accent}}{}{0em}{}
\titleformat{\subsection}{\normalfont\large\bfseries\color{accent}}{}{0em}{}
\titleformat{\subsubsection}{\normalfont\normalsize\bfseries\color{accent!85!black}}{}{0em}{}
\titlespacing*{\section}{0pt}{14pt}{6pt}
\titlespacing*{\subsection}{0pt}{10pt}{4pt}
\setcounter{tocdepth}{2}

% ---- En-tête / pied de page (institution · page · projet) -------------------
\usepackage{fancyhdr}
\usepackage{lastpage}
\pagestyle{fancy}
\fancyhf{}
\fancyfoot[L]{\footnotesize\color{accent}TP Cybersécurité — Julien}
\fancyfoot[C]{\footnotesize\thepage\,/\,\pageref{LastPage}}
\fancyfoot[R]{\footnotesize\color{accent}Durcissement Claude Code · Docker}
\renewcommand{\headrulewidth}{0pt}
\renewcommand{\footrulewidth}{0.4pt}
\renewcommand{\footrule}{\hbox to\headwidth{\color{rulecol}\leaders\hrule height \footrulewidth\hfill}}

% ---- Encadrés lavande (mappe les blockquotes markdown) ----------------------
\usepackage[framemethod=default]{mdframed}
\newmdenv[
  backgroundcolor=boxbg, linecolor=boxframe, linewidth=0.6pt,
  skipabove=7pt, skipbelow=7pt, leftmargin=0pt, rightmargin=0pt,
  innerleftmargin=9pt, innerrightmargin=9pt, innertopmargin=6pt,
  innerbottommargin=6pt, roundcorner=2pt
]{calloutbox}
\let\oldquote\quote \let\endoldquote\endquote
\renewenvironment{quote}{\begin{calloutbox}\small}{\end{calloutbox}}

% ---- Tables : booktabs + réduction pour tables larges -----------------------
\usepackage{booktabs}
\usepackage{longtable}
\usepackage{array}
\usepackage{etoolbox}
\setlength{\tabcolsep}{4pt}
\renewcommand{\arraystretch}{1.15}
\AtBeginEnvironment{longtable}{\footnotesize}
\AtBeginEnvironment{tabular}{\small}

% ---- Code : fvextra (retour à la ligne) + fond léger ------------------------
\usepackage{fvextra}
\fvset{breaklines=true, breakanywhere=true, fontsize=\small, commandchars=\\\{\}}

% ---- Liens : couleurs passées via pandoc -V (colorlinks/linkcolor/urlcolor)
%      -> appliquées dans le \hypersetup du template, où 'accent' est déjà défini.

% ---- Espacements généraux ---------------------------------------------------
\usepackage{parskip}       % paragraphes séparés par un blanc (pas d'indentation)
\usepackage{enumitem}
\setlist{itemsep=1pt, topsep=3pt, parsep=0pt}
\usepackage{graphicx}
\usepackage{adjustbox}   % réduit une figure trop large à \linewidth
\usepackage{tikz}
\usetikzlibrary{positioning, arrows.meta, fit, backgrounds, shapes.geometric, calc}
% ---- Styles TikZ (charte) ---------------------------------------------------
\tikzset{
  ebox/.style={draw=boxframe, fill=boxbg, rounded corners=2pt, align=center,
               font=\footnotesize, inner sep=5pt, line width=0.5pt},
  hbox/.style={draw=accent, fill=accent!8, rounded corners=2pt, align=center,
               font=\footnotesize, inner sep=5pt, line width=0.8pt},
  dbox/.style={draw=red!55!black, fill=red!5, rounded corners=2pt, align=center,
               font=\footnotesize, inner sep=5pt, line width=0.6pt},
  gbox/.style={draw=green!45!black, fill=green!7, rounded corners=2pt, align=center,
               font=\footnotesize, inner sep=5pt, line width=0.6pt},
  grp/.style={draw=accent!65, rounded corners=3pt, dashed, line width=0.6pt},
  grpr/.style={draw=red!45!black, rounded corners=3pt, line width=0.7pt},
  grpg/.style={draw=green!40!black, rounded corners=3pt, line width=0.7pt},
  fl/.style={-{Stealth[length=5pt]}, line width=0.6pt, font=\scriptsize},
  fld/.style={-{Stealth[length=5pt]}, line width=0.6pt, dashed, font=\scriptsize},
  lbl/.style={font=\scriptsize\bfseries, text=accent},
}
